\section*{Hoofdstuk \ref{ch:g6:inside-the-earth} --- Het binnenste van de Aarde}

\begin{solution}{exo:g6:inside-the-earth:1}
De \omterm{def:g6:inside-the-earth:layers}{aardkorst}, de \omterm{def:g6:inside-the-earth:layers}{mantel}, de \omterm{def:g6:inside-the-earth:layers}{buitenkern}, de \omterm{def:g6:inside-the-earth:layers}{binnenkern}. De \omterm{def:g6:inside-the-earth:layers}{buitenkern} is
vloeibaar. De buiten- en de \omterm{def:g6:inside-the-earth:layers}{binnenkern} zijn van metaal (vooral ijzer, met
wat nikkel); de \omterm{def:g6:inside-the-earth:layers}{aardkorst} en de \omterm{def:g6:inside-the-earth:layers}{mantel} zijn van \omterm{def:g3:stones-pebbles-sand:rock}{gesteente}.
\end{solution}

\begin{solution}{exo:g6:inside-the-earth:2}
(a) Niet waar: de \omterm{def:g6:inside-the-earth:layers}{aardkorst} is dikker onder de \omterm{def:g3:continents-and-oceans:continent}{continenten} (ongeveer
30~km) dan onder de \omterm{def:g2:water-on-earth:ocean}{oceanen} (ongeveer 5~km). (b) Niet waar: de \omterm{def:g6:inside-the-earth:layers}{kern} is van
metaal, een legering van ijzer en nikkel. (c) Waar. (d) Waar: 12\,262~m,
ongeveer 12.3~km.
\end{solution}

\begin{solution}{exo:g6:inside-the-earth:3}
(a) Op 2891~km, ongeveer 2900~km. (b) Op 5150~km. (c) $5150 - 2891 = 2259$:
ongeveer 2260~km.
\end{solution}

\begin{solution}{exo:g6:inside-the-earth:4}
``De schaal is als de \underline{\omterm{def:g6:inside-the-earth:layers}{aardkorst}}, het eiwit is als de
\underline{\omterm{def:g6:inside-the-earth:layers}{mantel}} en de dooier is als de \underline{\omterm{def:g6:inside-the-earth:layers}{kern}}. Maar anders dan
de dooier heeft de \underline{\omterm{def:g6:inside-the-earth:layers}{kern}} twee delen: een vloeibare
\underline{\omterm{def:g6:inside-the-earth:layers}{buitenkern}} en een vaste \underline{\omterm{def:g6:inside-the-earth:layers}{binnenkern}}.''
\end{solution}

\begin{solution}{exo:g6:inside-the-earth:5}
$8.849 \div 6371 \approx 0.00139$, en $0.00139 \times 100 \approx 0.14$:
Everest is maar ongeveer \textbf{0.14~\%} van de straal. Van ver gezien is
de Aarde een heel gladde bol: haar hoogste \omterm{def:g1:landscapes-around-us:landform}{berg} is een piepklein bultje.
\end{solution}

\begin{solution}{exo:g6:inside-the-earth:6}
(a) $30 \div 5 = 6$: zes keer dikker. (b) $100 \div 5 = 20$: twintig keer
dikker.
\end{solution}

\begin{solution}{exo:g6:inside-the-earth:7}
(a) $12 + 1 \times 25 = 37~^\circ$C. (b) $12 + 1.5 \times 25 = 12 + 37.5 =
49.5~^\circ$C. (c) $12 + 2 \times 25 = 62~^\circ$C.
\end{solution}

\begin{solution}{exo:g6:inside-the-earth:8}
De stijging is $74 - 14 = 60$ graden over 2~km, dus $60 \div 2 =
\textbf{30}$~graden per kilometer. Ja: 30 ligt tussen het gemiddelde van
het oosten (25) en dat van het westen (34).
\end{solution}

\begin{solution}{exo:g6:inside-the-earth:9}
(a) Vlak boven de \omterm{def:g6:inside-the-earth:layers}{kern} ongeveer 13.7~km per seconde; vlak onder de
bovenkant ervan ongeveer 8.1~km per seconde: de P-golven worden plotseling
langzamer. (b) Tussen 2891~km en 5150~km, in de \omterm{def:g6:inside-the-earth:layers}{buitenkern}. (c) S-golven
kunnen niet door een vloeistof reizen: de \omterm{def:g6:inside-the-earth:layers}{buitenkern} is vloeibaar.
\end{solution}

\begin{solution}{exo:g6:inside-the-earth:10}
(a) $12.3 \div 30 = 0.41$: ongeveer \textbf{41~\%} van een gemiddelde
continentale korst. (b) Nee: onder een \omterm{def:g3:continents-and-oceans:continent}{continent} begint de \omterm{def:g6:inside-the-earth:layers}{mantel} onder
de korst, gemiddeld ongeveer 30~km diep; het boorgat stopte ver daarboven,
in de korst.
\end{solution}

\begin{solution}{exo:g6:inside-the-earth:11}
(a) $29 \times 6371 = 184\,759$, en $15 + 184\,759 = 184\,774$: ongeveer
185\,000~graden. (b) Het laboratorium geeft ongeveer 6000~graden aan de
rand van de \omterm{def:g6:inside-the-earth:layers}{binnenkern}: het getal van Leo is ongeveer 30~keer te groot
($184\,774 \div 6000 \approx 31$). (c) Het tempo van 29~graden per
kilometer is alleen in de eerste 2~km of zo gemeten. Leo gebruikt het
duizenden kilometers voorbij de metingen; diep onder de grond stijgt de
temperatuur veel langzamer. Een tempo mag je alleen gebruiken waar het is
gemeten.
\end{solution}

\begin{solution}{exo:g6:inside-the-earth:12}
(a) Als de \omterm{def:g6:inside-the-earth:layers}{aardkorst} lichter is dan het gemiddelde, moet het materiaal
dieper in de Aarde dichter zijn dan 5.5, om het gemiddelde omhoog te
brengen. (b) In de \omterm{def:g6:inside-the-earth:layers}{buitenkern} en de \omterm{def:g6:inside-the-earth:layers}{binnenkern} (van ongeveer 9.9 tot
ongeveer 13~gram per kubieke centimeter); de \omterm{def:g6:inside-the-earth:layers}{mantel} haalt helemaal onderin
maar net 5.6. (c) Dichtere materialen zakken onder lichtere, zoals een
steen in water zinkt: het zware metaal is in het midden terechtgekomen en
de lichtere \omterm{def:g3:stones-pebbles-sand:rock}{gesteenten} bovenop.
\end{solution}

\begin{solution}{pb:g6:inside-the-earth:1}
\textbf{1.} $2 \times 6371 = 12\,742$~km. Het model is 2~m breed, dat is
2000~mm.

\textbf{2.} $12\,742 \div 2000 \approx 6.37$: 1~mm van het model staat voor
ongeveer 6.37~km van de echte Aarde.

\textbf{3.} $8.849 \div 6.371 \approx 1.39$: Everest is op het model
ongeveer \textbf{1.4~mm} hoog, een zandkorrel.

\textbf{4.} $12.3 \div 6.371 \approx 1.9$~mm. Het middelpunt van het model
ligt 1000~mm (1~m) onder het oppervlak: het boorgat is ongeveer 520~keer
korter ($1000 \div 1.9 \approx 520$).

\textbf{5.} $2891 \div 6.371 \approx 453.8$: de bovenkant van de \omterm{def:g6:inside-the-earth:layers}{buitenkern}
wordt ongeveer \textbf{454~mm} onder het oppervlak geschilderd.

\textbf{6.} $5150 \div 6.371 \approx 808.3$: ongeveer 808~mm onder het
oppervlak. De straal van de \omterm{def:g6:inside-the-earth:layers}{binnenkern} op het model is $1000 - 808 = 192$~mm.

\textbf{7.} De \omterm{def:g6:inside-the-earth:layers}{mantel}: van ongeveer 5~mm tot 454~mm, ongeveer 449~mm. De
\omterm{def:g6:inside-the-earth:layers}{buitenkern} neemt $808 - 454 = 354$~mm in en de \omterm{def:g6:inside-the-earth:layers}{binnenkern} 192~mm.

\textbf{8.} $30 \div 6.371 \approx 4.7$~mm: een dunne band, die met een
fijne kwast geschilderd kan worden.

\textbf{9.} $15 + 2 \times 29 = 73~^\circ$C. Op het model ligt 2~km
$2 \div 6.371 \approx 0.31$~mm onder het oppervlak: nauwelijks de dikte van
drie vellen papier.

\textbf{10.} $4 \div 6.371 \approx 0.63$~mm, tegen ongeveer 1.9~mm voor het
boorgat: de mijn is minder diep, met ongeveer $1.9 - 0.6 = 1.3$~mm.

\textbf{11.} Of een stof smelt, hangt af van de temperatuur en van de druk.
De \omterm{def:g6:inside-the-earth:layers}{binnenkern} wordt het hardst van allemaal samengeperst, met ongeveer
3.3~miljoen keer de luchtdruk: samengeperst ijzer smelt pas bij een hogere
temperatuur, dus blijft het ijzer van de \omterm{def:g6:inside-the-earth:layers}{binnenkern} vast, hoewel het heter
is dan de \omterm{def:g6:inside-the-earth:layers}{buitenkern}.

\textbf{12.} $5 \div 6.371 \approx 0.785$: de oceanische korst moet op het
model van 2~meter maar \textbf{0.78~mm} dik worden geschilderd, de dikte
van ongeveer acht vellen papier ($0.78 \div 0.1 \approx 8$).
\end{solution}
